%%% variant: default
\section{Smoke}

This is \textbf{bold} text.
Option checks: \cite{single} and \parencite{group}.
The inline delimiters $x^2$ and \(y^2\) both work.

\[
a = b
\]

\[
c = d
\]

\[
\begin{aligned}
z_1 &= w_{11}x_1+w_{12}x_2+b_1, \\
z_2 &= w_{21}x_1+w_{22}x_2+b_2.
\end{aligned}
\]

\[\begin{pmatrix}
u & v \\
w & z
\end{pmatrix}\]


\begin{equation}
e = mc^2
\end{equation}

\begin{quote}
\textbf{Beobachtungen} → Daten → erkennbare Zusammenhänge → Modell → Anwendung auf neue Fälle
\end{quote}

\par\addvspace{\medskipamount}
\noindent\begin{tabularx}{\linewidth}{@{}>{\raggedright\arraybackslash}X>{\raggedright\arraybackslash}X@{}}
\toprule
Symbol & Bedeutung \\
\midrule
$x$ & Eingabe \\
\bottomrule
\end{tabularx}
\par\addvspace{\medskipamount}

% This must not appear in the PDF: HIDDENCOMMENTMARKER 
%
%# Hidden heading
%\undefinedCommentCommand{ignored}
%
%TODO: complete later.
%
Visible before % inline comment 
and after.

%%% variant: instructor
\section{Smoke}

This is \textbf{bold} text.
Option checks: \cite{single} and \parencite{group}.
The inline delimiters $x^2$ and \(y^2\) both work.

\[
a = b
\]

\[
c = d
\]

\[
\begin{aligned}
z_1 &= w_{11}x_1+w_{12}x_2+b_1, \\
z_2 &= w_{21}x_1+w_{22}x_2+b_2.
\end{aligned}
\]

\[\begin{pmatrix}
u & v \\
w & z
\end{pmatrix}\]


\begin{equation}
e = mc^2
\end{equation}

\begin{quote}
\textbf{Beobachtungen} → Daten → erkennbare Zusammenhänge → Modell → Anwendung auf neue Fälle
\end{quote}

\par\addvspace{\medskipamount}
\noindent\begin{tabularx}{\linewidth}{@{}>{\raggedright\arraybackslash}X>{\raggedright\arraybackslash}X@{}}
\toprule
Symbol & Bedeutung \\
\midrule
$x$ & Eingabe \\
\bottomrule
\end{tabularx}
\par\addvspace{\medskipamount}

% This must not appear in the PDF: HIDDENCOMMENTMARKER 
%
%# Hidden heading
%\undefinedCommentCommand{ignored}
%
%TODO: complete later.
%
Visible before % inline comment 
and after.

%%% variant: handout
\section{Smoke}

This is \textbf{bold} text.
Option checks: \cite{single} and \parencite{group}.
The inline delimiters $x^2$ and \(y^2\) both work.

\[
a = b
\]

\[
c = d
\]

\[
\begin{aligned}
z_1 &= w_{11}x_1+w_{12}x_2+b_1, \\
z_2 &= w_{21}x_1+w_{22}x_2+b_2.
\end{aligned}
\]

\[\begin{pmatrix}
u & v \\
w & z
\end{pmatrix}\]


\begin{equation}
e = mc^2
\end{equation}

\begin{quote}
\textbf{Beobachtungen} → Daten → erkennbare Zusammenhänge → Modell → Anwendung auf neue Fälle
\end{quote}

\par\addvspace{\medskipamount}
\noindent\begin{tabularx}{\linewidth}{@{}>{\raggedright\arraybackslash}X>{\raggedright\arraybackslash}X@{}}
\toprule
Symbol & Bedeutung \\
\midrule
$x$ & Eingabe \\
\bottomrule
\end{tabularx}
\par\addvspace{\medskipamount}

% This must not appear in the PDF: HIDDENCOMMENTMARKER 
%
%# Hidden heading
%\undefinedCommentCommand{ignored}
%
%TODO: complete later.
%
Visible before % inline comment 
and after.

%%% variant: slides
\begin{frame}[t]
\frametitle{Smoke}

\end{frame}

\begin{frame}[t]

This is \textbf{bold} text.
Option checks: \cite{single} and \parencite{group}.
The inline delimiters $x^2$ and \(y^2\) both work.

\[
a = b
\]

\[
c = d
\]

\[
\begin{aligned}
z_1 &= w_{11}x_1+w_{12}x_2+b_1, \\
z_2 &= w_{21}x_1+w_{22}x_2+b_2.
\end{aligned}
\]

\[\begin{pmatrix}
u & v \\
w & z
\end{pmatrix}\]


\begin{equation}
e = mc^2
\end{equation}

\begin{quote}
\textbf{Beobachtungen} → Daten → erkennbare Zusammenhänge → Modell → Anwendung auf neue Fälle
\end{quote}

\par\addvspace{\medskipamount}
\noindent\begin{tabularx}{\linewidth}{@{}>{\raggedright\arraybackslash}X>{\raggedright\arraybackslash}X@{}}
\toprule
Symbol & Bedeutung \\
\midrule
$x$ & Eingabe \\
\bottomrule
\end{tabularx}
\par\addvspace{\medskipamount}

% This must not appear in the PDF: HIDDENCOMMENTMARKER 
%
%# Hidden heading
%\undefinedCommentCommand{ignored}
%
%TODO: complete later.
%
Visible before % inline comment 
and after.

\end{frame}

%%% variant: slides-dozent-h1-frame
\begin{frame}[t]
\frametitle{Smoke}

This is \textbf{bold} text.
Option checks: \cite{single} and \parencite{group}.
The inline delimiters $x^2$ and \(y^2\) both work.

\[
a = b
\]

\[
c = d
\]

\[
\begin{aligned}
z_1 &= w_{11}x_1+w_{12}x_2+b_1, \\
z_2 &= w_{21}x_1+w_{22}x_2+b_2.
\end{aligned}
\]

\[\begin{pmatrix}
u & v \\
w & z
\end{pmatrix}\]


\begin{equation}
e = mc^2
\end{equation}

\begin{quote}
\textbf{Beobachtungen} → Daten → erkennbare Zusammenhänge → Modell → Anwendung auf neue Fälle
\end{quote}

\par\addvspace{\medskipamount}
\noindent\begin{tabularx}{\linewidth}{@{}>{\raggedright\arraybackslash}X>{\raggedright\arraybackslash}X@{}}
\toprule
Symbol & Bedeutung \\
\midrule
$x$ & Eingabe \\
\bottomrule
\end{tabularx}
\par\addvspace{\medskipamount}

% This must not appear in the PDF: HIDDENCOMMENTMARKER 
%
%# Hidden heading
%\undefinedCommentCommand{ignored}
%
%TODO: complete later.
%
Visible before % inline comment 
and after.

\end{frame}

%%% variant: slides-h1-hide-level-1
\begin{frame}[t]

This is \textbf{bold} text.
Option checks: \cite{single} and \parencite{group}.
The inline delimiters $x^2$ and \(y^2\) both work.

\[
a = b
\]

\[
c = d
\]

\[
\begin{aligned}
z_1 &= w_{11}x_1+w_{12}x_2+b_1, \\
z_2 &= w_{21}x_1+w_{22}x_2+b_2.
\end{aligned}
\]

\[\begin{pmatrix}
u & v \\
w & z
\end{pmatrix}\]


\begin{equation}
e = mc^2
\end{equation}

\begin{quote}
\textbf{Beobachtungen} → Daten → erkennbare Zusammenhänge → Modell → Anwendung auf neue Fälle
\end{quote}

\par\addvspace{\medskipamount}
\noindent\begin{tabularx}{\linewidth}{@{}>{\raggedright\arraybackslash}X>{\raggedright\arraybackslash}X@{}}
\toprule
Symbol & Bedeutung \\
\midrule
$x$ & Eingabe \\
\bottomrule
\end{tabularx}
\par\addvspace{\medskipamount}

% This must not appear in the PDF: HIDDENCOMMENTMARKER 
%
%# Hidden heading
%\undefinedCommentCommand{ignored}
%
%TODO: complete later.
%
Visible before % inline comment 
and after.

\end{frame}

